\noindent\textbf{Por qué se usa este procedimiento.} La ley de giro es $\theta=TL/(GJ)$. Despejar $T=GJ\theta/L$ exige convertir el grado a $\pi/180$ rad. En un tubo, $J=\pi(R_o^4-R_i^4)/2$ utiliza radios, no diámetros.

\begin{enumerate}

\item Un grado son $\pi/180$ rad y $J = \pi(25^4 - 20^4)/2$ mm$^4$.

\end{enumerate}

\noindent\textbf{Operaciones y resultados intermedios.}

\begin{enumerate}

\item Torsor máximo en N mm.

\[\fitmath{y_{1}=\frac{26950\cdot \frac{\pi\cdot \left({25}^{4}-{20}^{4}\right)}{2}\cdot \frac{\pi}{180}}{1000}\approx 170397.2\,\mathrm{N mm}}\]

\item Torsor máximo en N m.

\[\fitmath{y_{2}=\frac{\frac{26950\cdot \frac{\pi\cdot \left({25}^{4}-{20}^{4}\right)}{2}\cdot \frac{\pi}{180}}{1000}}{1000}\approx 170.3972\,\mathrm{N m}}\]

\end{enumerate}
